\section{Adaptive channels on arbitrary inputs}

Finite local measurements and classical feedforward act on mixed states by
summing unnormalized completely positive operations. This includes the
finite-round extension discussed in~\cite[paragraph ``State transformations
    with QC and LOCC'']{Piroli2021LOCC}. The number of rounds is not bounded
independently of system size here; this condition is separate from circuit
depth in the source's asymptotic phase relation, as explained
in~\cite{gap:psc21_adaptive_channel_round_scope}.

\begin{definition}[Onsite instruments with conditional continuation]
    \label{def:qc_onsite_feedforward}
    \uses{def:qc_onsite_channel}
    \lean{QuantumCircuit.OnsiteChannel.feedforwardMap}
    \leanok
    An onsite instrument has local Kraus operators $K_{i,j}$ satisfying
    $\sum_jK_{i,j}^{\dagger}K_{i,j}=I$. For the communicated tuple of outcomes
    $J$, let $K_J=\bigotimes_iK_{i,J(i)}$ and let $\Psi_J$ be the selected
    continuation. Averaging the outcomes gives the map
    \begin{align}
        \mathcal F(\rho) & =\sum_J\Psi_J(K_J\rho K_J^{\dagger}).
        \label{eq:qc_onsite_feedforward}
    \end{align}
    Each summand uses an unnormalized post-measurement operator.
\end{definition}

\begin{theorem}[Trace preservation with classical feedforward]
    \label{thm:qc_onsite_feedforward_cptp}
    \uses{def:qc_onsite_feedforward}
    \lean{QuantumCircuit.OnsiteChannel.feedforwardMap_isKrausCPTP,
        QuantumCircuit.OnsiteChannel.feedforwardMap_const,
        QuantumCircuit.OnsiteChannel.comp_feedforwardMap}
    \leanok
    If each $\Psi_J$ is trace preserving and completely positive, then so is
    $\mathcal F$. If $\Psi_J=\Psi$ is independent of $J$, then
    $\mathcal F=\Psi\circ\Phi$, where $\Phi$ is the channel of the onsite
    instrument with its outcomes forgotten. A further linear map $\Theta$
    may be composed into every continuation in
    (\ref{eq:qc_onsite_feedforward}).
\end{theorem}
\begin{proof}
    \leanok
    Let $B_{J,a}$ be Kraus operators for $\Psi_J$. Then $B_{J,a}K_J$ are
    Kraus operators for $\mathcal F$, and
    \begin{align}
        \sum_{J,a}(B_{J,a}K_J)^{\dagger}(B_{J,a}K_J)
         & =\sum_JK_J^{\dagger}\Bigl(\sum_aB_{J,a}^{\dagger}B_{J,a}\Bigr)K_J
        =\sum_JK_J^{\dagger}K_J=I. \notag
    \end{align}
    The two composition identities follow by linearity of the finite sum.
\end{proof}

\begin{definition}[Finite adaptive local protocol]
    \label{def:qc_adaptive_protocol}
    \uses{def:qc_onsite_feedforward, def:qc_local_channel_circuit}
    \lean{QuantumCircuit.IsAdaptiveChannelProtocol,
        QuantumCircuit.IsAdaptiveChannelConversion}
    \leanok
    A finite adaptive local protocol is a tree formed from onsite channels,
    nearest-neighbor channel layers, and onsite instruments with a
    continuation chosen according to the complete communicated outcome tuple.
    A depth bound $T$ bounds the number of nearest-neighbor layers on every
    branch. Onsite channels and instruments have zero cost, and one layer
    adds one to the depth. A larger depth bound also bounds the same protocol.
    An operator $\rho$ is adaptively converted into $\sigma$ in depth $T$
    when $\Psi(\rho)=\sigma$ for such a protocol $\Psi$.
\end{definition}

\begin{theorem}[Channels, composition, and deterministic inclusion]
    \label{thm:qc_adaptive_protocol_properties}
    \uses{def:qc_adaptive_protocol, thm:qc_onsite_feedforward_cptp}
    \lean{QuantumCircuit.IsAdaptiveChannelProtocol.isKrausCPTP,
        QuantumCircuit.IsAdaptiveChannelProtocol.onsite_comp,
        QuantumCircuit.IsAdaptiveChannelProtocol.comp,
        QuantumCircuit.IsAdaptiveChannelProtocol.channelLayer,
        QuantumCircuit.IsLocalChannelProtocol.adaptive,
        QuantumCircuit.IsAdaptiveChannelConversion.refl,
        QuantumCircuit.IsAdaptiveChannelConversion.mono,
        QuantumCircuit.IsAdaptiveChannelConversion.trans,
        QuantumCircuit.IsAdaptiveChannelConversion.density,
        QuantumCircuit.IsLocalChannelConversion.adaptive}
    \leanok
    Every finite adaptive protocol is trace preserving and completely
    positive. Two protocols with depth bounds $T_1,T_2$ compose with depth
    bound $T_1+T_2$. Deterministic local channel protocols are adaptive
    protocols with the same depth bound. Adaptive conversion is reflexive
    in depth zero, monotone in its depth bound, transitive with additive
    depth, and preserves density matrices.
\end{theorem}
\begin{proof}
    \leanok
    Induct on the protocol tree. The onsite and layer cases use their
    channel properties, and the instrument case uses
    Theorem~\ref{thm:qc_onsite_feedforward_cptp}. To compose protocols, put
    the second protocol after each branch of the first. An instrument keeps
    the same branch depth, a nearest-neighbor layer adds one, and enlarging
    a depth bound commutes with adding the second depth. An onsite operation
    followed by a fixed continuation is an instrument whose continuation
    ignores its outcomes. This also embeds deterministic protocols.
    The remaining conversion properties follow from the identity channel,
    the depth bounds, composition, and preservation of positivity and trace.
\end{proof}

\begin{definition}[The average channel of a measurement round]
    \label{def:qc_measurement_round_channel}
    \uses{def:qc_measurement_rounds}
    \lean{QuantumCircuit.MeasurementRound.map}
    \leanok
    For a measurement round with circuit $U$, outcome projectors $P_m$, and
    conditional product unitaries $V_m$, its average channel is
    $\rho\mapsto\sum_m V_mP_mU\rho U^{\dagger}P_mV_m^{\dagger}$.
\end{definition}

\begin{theorem}[Measurement rounds act on density matrices]
    \label{thm:qc_measurement_round_channel}
    \uses{def:qc_measurement_round_channel, def:qc_measurement_rounds}
    \lean{QuantumCircuit.MeasurementRound.sum_kraus,
        QuantumCircuit.MeasurementRound.map_isKrausCPTP,
        QuantumCircuit.MeasurementRound.map_vecMulVec,
        QuantumCircuit.MeasurementRound.map_vecMulVec_of_isPreparationOf}
    \leanok
    The average map of a measurement round is trace preserving and
    completely positive. On a pure input $\ket v\bra v$, its output is
    $\sum_m\ket{V_mP_mUv}\bra{V_mP_mUv}$. In particular, if the round
    deterministically prepares $\psi$ from a product input $v$, and both
    input and target pure operators have trace one, its average output is
    $\ket\psi\bra\psi$.
\end{theorem}
\begin{proof}
    \leanok
    Unitarity of $U,V_m$, orthogonality of each $P_m$, and
    $\sum_mP_m=I$ give
    $\sum_m(V_mP_mU)^{\dagger}(V_mP_mU)=I$.
    Conjugating a rank-one operator gives the stated pure-input formula.
    For deterministic preparation, every output is $c_m\psi$, including
    $c_m=0$ for an outcome of probability zero. Thus the average output is
    $(\sum_m|c_m|^2)\ket\psi\bra\psi$. Trace preservation and normalization
    give $\sum_m|c_m|^2=1$.
\end{proof}
